\section{Introduction}
Let $(M,\omega)$ be a compact connected symplectic manifold equipped with an action of a compact connected Lie group $G$ by symplectomorphisms. Suppose that the action of $G$ is Hamiltonian, meaning that there is a $G$-equivariant map, the moment map,
\[ \mu_\g \colon \ M \rightarrow \g^\ast,\]
where $\g^\ast$ is the dual of the Lie algebra $\g=\tn{Lie}(G)$, satisfying the moment map condition
\begin{gather}\label{eqn:mm}
\iota(X_M)\omega=-{\rm d}\pair{\mu_\g}{X}, \qquad X \in \g.
\end{gather}
Let $\big(L,\nabla^L\big)$ be a $G$-equivariant prequantum line bundle with connection on $M$, i.e., $L$~is a~$G$-equivariant Hermitian line bundle with compatible connection $\nabla^L$, $\big(\nabla^L\big)^2=-2\pi \i \omega$ and the derivative of the $G$-action on $L$ satisfies Kostant's condition
\[ \L^L_X-\nabla^L_{X_M}=2\pi \i\pair{\mu_\g}{X}.\]
Choose a compatible almost complex structure $J$ on $M$, i.e., $\omega(Jw,Jv)=\omega(w,v)$ and $\omega(w,Jv)\allowbreak =:g(w,v)$ is a Riemannian metric. Let $D_L$ denote the Dolbeault--Dirac operator twisted by $\big(L,\nabla^L\big)$, an elliptic differential operator acting on sections of the spinor bundle $\wedge T^\ast_{0,1}M\otimes L$. The kernel of~$D_L$ carries an action of~$G$, and the $G$-equivariant index is defined to be the difference $\index_G(D_L):=\ker\big(D_L^{\tn{even}}\big)-\ker\big(D_L^{\tn{odd}}\big)$ of the kernel of~$D_L$ on even/odd degree forms, regarded as an element of the representation ring~$R(G)$.

The quantization-commutes-with-reduction theorem ($[Q,R]=0$ theorem) describes the multiplicity of the trivial representation in $\index_G(D_L)$ in terms of the symplectic quotient $M^{\tn{red}}:=\mu_\g^{-1}(0)/G$. When $0$ is a regular value of $\mu_\g$, $M^{\tn{red}}$ is an orbifold and the theorem states that $\index_G(D)^G$ equals the index of the twisted Dolbeault--Dirac operator $D^{\tn{red}}_{L^{\tn{red}}}$ on $M^{\tn{red}}$. The theorem was first conjectured by Guillemin--Sternberg~\cite{GuilleminSternbergConjecture}, and the general case ($M$, $G$ both compact, $0$ a regular value) was first proved by Meinrenken~\cite{MeinrenkenSymplecticSurgery}. Different proofs of the $[Q,R]=0$ theorem were given by Tian--Zhang~\cite{TianZhang} and Paradan~\cite{ParadanRiemannRoch}. The theorem has since been extended in various directions.

There are versions of the $[Q,R]=0$ theorem when $0$ is not necessarily a regular value, due to Meinrenken--Sjamaar~\cite{MeinrenkenSjamaar}; below we will give a precise statement of one of these results, involving a partial \emph{shift desingularization}, i.e., $\index_G(D_L)^G$ is related to the index on the symplectic quotient at a nearby weakly regular value. At the same time, we introduce some notation that will be of use later on.

Fix a maximal torus $T$ with Lie algebra $\t$. Let $\Lambda \subset \t^\ast$ be the (real) weight lattice. Given $\lambda \in \Lambda$, the corresponding character $T \rightarrow U(1)$ is written $t \mapsto t^\lambda={\rm e}^{2\pi \i\pair{\lambda}{X}}$ where $t={\rm e}^X$, $X \in \t$. Let $\R \subset \Lambda$ be the set of roots. We also fix a closed positive Weyl chamber $\t_+$, which determines a set of positive (resp.\ negative) roots $\R_{\pm}$. For each relatively open face $\sigma \subset \t_+^\ast$, the stabilizer~$G_\xi$ of points~$\xi \in \sigma$ under the coadjoint action, does not depend on~$\xi$, and will be denoted $G_\sigma$. If $\sigma_1 \subset \ol{\sigma}_2$ then $G_{\sigma_1}\supset G_{\sigma_2}$. Note also that $G_\sigma$ is connected and contains the maximal torus~$T$. The Lie algebra~$\g_\sigma$ decomposes into its semi-simple and central parts
$ \g_\sigma=[\g_\sigma,\g_\sigma]\oplus \z_\sigma$.
The subspace $\z_\sigma^\ast \subset \t^\ast$ is defined to be the annihilator of $[\g_\sigma,\g_\sigma]$, or equivalently the fixed point set of the coadjoint $G_\sigma$ action. The face $\sigma$ is an open subset of $\z_\sigma^\ast$.

Let $\Delta=\mu_\g(M)\cap \t_+^\ast$ be the moment polytope. A well-known theorem in symplectic geometry states that there is a~unique face $\sigma \subset \t^\ast_+$ of minimal dimension such that $\Delta \subset \ol{\sigma}$ (briefly, this is a~consequence of~\eqref{eqn:mm}, which implies that ${\rm d}\mu_\g$ has constant rank on the top dimensional $G$-orbit type stratum, and the complement of the latter has codimension at least~$2$); $\sigma$ is called the \emph{principal face} or \emph{principal wall}. The corresponding symplectic cross-section, called the \emph{principal cross-section},
$Y=\mu_\g^{-1}(\sigma)$
is a Hamiltonian $G_\sigma$-space. Moreover the semi-simple part $[G_\sigma,G_\sigma]$ of $G_\sigma$ acts trivially on $Y$. For further details, see for example~\cite{ConvexitySympCuts} and references therein.

Let $I \subset \z_\sigma^\ast$ be the smallest affine subspace containing $\Delta$. Let $\t_I\subset \t$ be the annihilator of the subspace parallel to $I$, and let $T_I=\exp(\t_I)\subset T$ be the corresponding subtorus. By equation~\eqref{eqn:mm}, $\t_I$ is the generic infinitesimal stabilizer of $Y$. In particular $T_I$ acts trivially, hence the quotient torus $T/T_I$ acts on $Y$. The moment map $\mu_\g$ may have no non-trivial regular values. But the restriction
\[\mu_\g|_Y\colon \ Y \rightarrow I\]
viewed as a map with codomain $I$, always has non-trivial regular values, and we will refer to these as \emph{weakly-regular values}. If $\xi$ is a weakly-regular value, then the reduced space $M_\xi=\mu_\g^{-1}(\xi)/G_\sigma$ is an orbifold. Let $L_\xi=L|_{\mu_\g^{-1}(\xi)}/G_\sigma$ be the corresponding (orbifold) line bundle over $M_\xi$.

\begin{Theorem}[\cite{MeinrenkenSjamaar}, see also \cite{ParadanRiemannRoch,WittenNonAbelian}]
\label{thm:MeinSja}
Let $(M,\omega,\mu_\g)$ be a compact connected Hamiltonian $G$-space with moment polytope $\Delta$. If $0 \notin \Delta$ then $\index_G(D_L)^G=0$. Otherwise for every weakly-regular value $\xi \in \Delta$ sufficiently close to $0$, $\index_G(D_L)^G$ equals the index of the Dolbeault--Dirac operator $D^{\tn{red}}_{L_\xi}$ on the reduced space $M_\xi$.
\end{Theorem}

We will now describe the main result of this article and its relation to Theorem \ref{thm:MeinSja}. Consider tensor powers $L^k$, $k \in \bZ_{>0}$ of the prequantum line bundle. For a dominant weight $\lambda$, let $\chi_\lambda \in R(G)$ denote the character of the irreducible representation of $G$ with highest weight $\lambda$. We define the \emph{multiplicity function} $m_G(k,\lambda)$ by the expression
\begin{gather}
\label{eqn:multfcn}
\index_G(D_{L^k})=\sum_{\lambda \in \Lambda \cap \t_+^\ast} m_G(k,\lambda)\chi_\lambda.
\end{gather}
An important theme in the work of Szenes--Vergne \cite{SzenesVergne2010} and also in our approach, is that the function $m_G(k,\lambda)$ has more coherent behavior than its restriction to any fixed value of $k$.

The statement of the result requires some further background on orbifolds, for which we refer the reader to, for example, \cite[Appendix A]{daSilvaThesis}, \cite[Section 2]{MeinrenkenSymplecticSurgery}. A small warning is that we will not require the action of isotropy groups in orbifold charts to be effective (this is in agreement with the references \cite{daSilvaThesis,MeinrenkenSymplecticSurgery} mentioned above). One advantage of permitting this, is that for a locally free action of a compact Lie group $K$ on a manifold $P$, the corresponding orbifold $P/K$ has orbifold charts given automatically by the slice theorem, with the isotropy groups being simply the isotropy groups for the action of $K$ on $P$.

In fact all the orbifolds that we will encounter arise naturally as such quotients $P/K$, and one could avoid mentioning orbifolds altogether by working instead with suitable $K$-basic structures on $P$. An example is the description of characteristic forms for orbifold vector bundles, which can be defined in terms of orbifold charts for~$P/K$, or alternatively in terms of $K$-basic differential forms on $P$. In brief, the latter approach goes as follows. One can take the complex $(\Omega_{\tn{bas}}(P),{\rm d})$ of $K$-basic differential forms on $P$ as a working definition of the de Rham complex of~$P/K$ (if~$K$ acts freely then $P/K$ is a manifold and pullback of forms from $P/K$ to $P$ is an isomorphism of complexes $(\Omega(P/K),{\rm d})\simeq (\Omega_{\tn{bas}}(P),{\rm d})$). A~$K$-equivariant vector bundle $E \rightarrow P$ determines an orbifold vector bundle $E/K$ over $P/K$. Let $\theta$ be a connection on $P$ with curvature $F_\theta$. The choice of connection determines a \emph{Cartan map} (cf.~\cite{MeinrenkenEncyclopedia}) from closed $K$-equivariant forms $\alpha(X)$ on $P$ to closed $K$-basic forms: $\alpha(X) \mapsto \tn{Car}_\theta(\alpha):=\Pi_{\tn{hor}}\alpha(F_\theta)$, where $\Pi_{\tn{hor}}$ is the projection onto the horizontal part relative to the connection. The Cartan map induces an isomorphism from the $K$-equivariant cohomology of $P$ to the cohomology of the complex of basic differential forms on~$P$. If $\alpha(X)$ is a~$K$-equivariant characteristic form (constructed via the $K$-equivariant analogue of the usual Chern--Weil construction cf.~\cite{BerlineGetzlerVergne,MeinrenkenEncyclopedia}), then one may take $\tn{Car}_\theta(\alpha) \in \Omega_{\tn{bas}}(P)$ as the definition of the corresponding characteristic form for $E/K$.

Let $\xi \in \Delta$ be a weakly-regular value. By the moment map equation \eqref{eqn:mm}, the action of $K=T/T_I$ on the level set
\[ P=\mu_\g^{-1}(\xi) \]
is locally free. The set $S_P$ of elements $g \in T/T_I$ such that $P^g\ne \varnothing$ is finite. For each $g \in S_P$, we obtain an orbifold
\[ \Sigma_g=P^g/(T/T_I), \qquad \Sigma=\bigsqcup_{g \in S_P} \Sigma_g.\]
Note that $\Sigma_1=P/(T/T_I)=M_\xi$ identifies with the reduced space itself, and more generally $\Sigma_g$ identifies with a symplectic quotient of $Y^g$. For each $g \in S_P$ there is an immersion $\Sigma_g\hookrightarrow \Sigma$ induced by $P^g\hookrightarrow P$. Let $\nu_{\Sigma_g,\Sigma}$ denote the (orbifold) normal bundle (the quotient $\nu_{P^g,P}/(T/T_I)$), which inherits a complex structure from the almost complex structures on $Y$, $Y^g$. Define the characteristic form
\[ \calDC^g(\nu_{\Sigma_g,\Sigma})=\tn{det}_\bC\big(1-g_\nu^{-1} {\rm e}^{-\frac{\i}{2\pi}F_\nu}\big), \]
where $g_\nu$ denotes the action of $g$ on the normal bundle (defined in terms of an orbifold chart, or in terms of $\nu_{P^g,P}$), and $F_\nu$ denotes the curvature. Taking the quotient of $L|_{P^g}$ we obtain (orbifold) line bundles
\[ L_{\Sigma_g}=(L|_{P^g})/(T/T_I), \qquad L_\Sigma=\bigsqcup_{g \in S_P} L_{\Sigma_g}. \]
There is a locally constant function
\[ g_L \colon \ \Sigma_g \rightarrow U(1) \]
giving the phase of the action of $g$ on $L_{\Sigma_g}$ (or equivalently on $L|_{P^g}$). Let $d \colon \Sigma \rightarrow \bZ$ be the locally constant function giving the size of a generic isotropy group for $\Sigma$ (or equivalently the number of elements in the generic stabilizer for the $T/T_I$ action on $\sqcup P^g$).

Let $\theta$ be a connection for the locally free $K=T/T_I$-action on $\sqcup_{g \in S_P}P^g$. The curvature $F_\theta$ is horizontal and $\t/\t_I$-valued, hence for any $\lambda \in (\t/\t_I)^\ast=I$, the form $\pair{\lambda}{F_\theta}$ is $K$-basic, hence descends to $\Sigma$. With the preparations above, we can state the main result of this note.
\begin{Theorem}\label{thm:main}If $0 \notin \Delta$ then $m(k,0)=0$ for all $k \ge 1$. If $0 \in \Delta$ then there is a closed polytope $\p \subset \Delta$ of the same dimension as $\Delta$ and containing the origin such that the following is true. Let $C_\p$ denote the cone
\[ C_\p=\{(t,t\tau)\,|\,t \in (0,\infty), \tau \in \p \} \subset \bR \times \t^\ast.\]
Fix a weakly regular value $\xi \in \Delta$ sufficiently close to $0$ as in Theorem~{\rm \ref{thm:MeinSja}}. Let $P=\mu_\g^{-1}(\xi)$ and define $\Sigma$, $L_\Sigma$, etc.\ as above. Then for all $(k,\lambda) \in (\bZ_{>0}\times \Lambda)\cap C_\p$,
\begin{gather}\label{eqn:main}
m_G(k,\lambda)=\sum_{g \in S_P} g^{-\lambda} \int_{\Sigma_g} \frac{1}{d}\frac{g_L^k\Ch(L_\Sigma)^k\Td(\Sigma)}{\calDC^g(\nu_{\Sigma_g,\Sigma})}{\rm e}^{\pair{\lambda}{F_\theta}}.
\end{gather}
\end{Theorem}
